\begin{trapbox}
\textbf{一、夹角正确，未必是旋转；公式正确，也要用对中心。}
\begin{center}\begin{tikzpicture}[x=1cm,y=1cm,>=Stealth,font=\small]
\path[use as bounding box] (-6.6,-0.6) rectangle (6.6,5.9);

\node at (-3.3,5.4) {夹角相同，还要检查半径};
\begin{scope}[shift={(-4.8,1.3)},scale=1.15]
\draw[->,gray] (-0.3,0)--(2.8,0); \draw[->,gray] (0,-0.3)--(0,2.4);
\draw[blue!75!black,thick] (0,0)--(2,0) node[right] {$P$};
\draw[teal!75!black,thick] (0,0)--(45:2);
\draw[orange!85!black,dashed,thick] (0,0)--(45:1);
\fill[teal!75!black] (45:2) circle(2pt) node[above right] {$Q_{\text{对}}$};
\fill[orange!85!black] (45:1) circle(2pt) node[above left] {$Q_{\text{错}}$};
\draw[->] (0.5,0) arc(0:45:0.5);
\node[below left] at (0,0) {$B=O$};
\end{scope}
\node[align=center] at (-3.3,0.2) {$BP=BQ_{\text{对}}=2$\\$BQ_{\text{错}}=1$，角度却同为 $45^\circ$};
\node at (3.25,5.4) {中心不同，结果就不同};
\begin{scope}[shift={(3,2.2)},scale=0.48]
\draw[->,gray] (-4,0)--(6.4,0); \draw[->,gray] (0,-2)--(0,6.5);
\draw[blue!75!black,thick] (2,-1)--(5,3);
\draw[teal!75!black,very thick,->] (2,-1)--(-2,2);
\draw[orange!85!black,dashed,thick,->] (0,0)--(-3,5);
\draw[blue!75!black,dotted] (0,0)--(5,3);
\fill (2,-1) circle(3pt) node[below right] {$B(2,-1)$};
\fill (0,0) circle(3pt) node[below left] {$O$};
\fill[blue!75!black] (5,3) circle(3pt) node[right] {$P$};
\fill[teal!75!black] (-2,2) circle(3pt) node[left] {$Q_{\text{对}}$};
\fill[orange!85!black] (-3,5) circle(3pt) node[above] {$Q_{\text{错}}$};
\end{scope}
\node[align=center] at (3.25,0.2) {同为逆时针 $90^\circ$\\绕 $B$ 得 $(-2,2)$；绕 $O$ 得 $(-3,5)$};

\end{tikzpicture}\end{center}

\textbf{左图：} 错把 $Q_{\text{错}}$ 当成旋转结果，它与 $BP$ 夹角虽为 $45^\circ$，
却只有一半长度。应同时检查等长；若还存在长度比 $k$，应给旋转后的相对坐标乘 $k$。

\textbf{右图：} 错把 $(5,3)$ 直接绕原点转，得到 $(-3,5)$；
正确步骤是 $(5,3)-(2,-1)=(3,4)\to(-4,3)$，再加回 $(2,-1)$ 得 $(-2,2)$。
图中实线与虚线的起点不同，正是错误根源。
\end{trapbox}
\clearpage
\begin{trapbox}
\textbf{二、角度和转向各要核对，不能只背“交换坐标”。}
\begin{center}\begin{tikzpicture}[x=1cm,y=1cm,>=Stealth,font=\small]
\path[use as bounding box] (-6.6,-0.6) rectangle (6.6,5.6);

\node at (-3.3,5.1) {$45^\circ$ 与 $90^\circ$ 不能混用};
\begin{scope}[shift={(-4.1,0.9)},scale=0.82]
\draw[->,gray] (-1.7,0)--(3.7,0); \draw[->,gray] (0,-0.4)--(0,3.8);
\draw[blue!75!black,thick,->] (0,0)--(3,1) node[right] {$P(3,1)$};
\draw[teal!75!black,thick,->] (0,0)--({sqrt(2)},{2*sqrt(2)}) node[above right] {$Q_{45}$};
\draw[purple,dashed,thick,->] (0,0)--(-1,3) node[above left] {$Q_{90}$};
\draw[->,teal!75!black] (18.435:1) arc(18.435:63.435:1);
\draw[->,purple] (18.435:1.5) arc(18.435:108.435:1.5);
\node[below left] at (0,0) {$O$};
\end{scope}
\node at (-3.3,0) {$Q_{45}=(\sqrt2,2\sqrt2),\ Q_{90}=(-1,3)$};
\node at (3.3,5.1) {相同余弦，可以有相反转向};
\begin{scope}[shift={(2.4,2.5)},scale=1.8]
\draw[->,gray] (-0.4,0)--(1.9,0);
\draw[->,gray] (0,-0.9)--(0,1.2);
\draw[blue!75!black,thick,->] (0,0)--(1,0) node[right] {$P$};
\draw[teal!75!black,thick,->] (0,0)--(45:1) node[above right] {$Q_+$};
\draw[orange!85!black,dashed,thick,->] (0,0)--(-45:1) node[below right] {$Q_-$};
\draw[->,teal!75!black] (0.5,0) arc(0:45:0.5);
\draw[->,orange!85!black] (0.5,0) arc(0:-45:0.5);
\node[below left] at (0,0) {$B=O$};
\end{scope}
\node[align=center] at (3.3,0.1) {$\cos(\pm45^\circ)=\sqrt2/2$\\$\sin45^\circ>0,\quad\sin(-45^\circ)<0$};

\end{tikzpicture}\end{center}

\textbf{左图：} 把 $(3,1)$ 变为 $(-1,3)$ 是转 $90^\circ$；
转 $45^\circ$ 应用坐标和差除以 $\sqrt2$，得到 $(\sqrt2,2\sqrt2)$。

\textbf{右图：} 只知道余弦无法区分 $Q_+,Q_-$；只给普通夹角而不指定方向时，
两者都可能。顺时针 $45^\circ$ 应代 $-45^\circ$，得到下方的点。
反求角度时须同时检查正弦与余弦，并指定角度范围。

\textbf{不要把“正弦为负”直接等同于所有顺时针角。}
例如逆时针 $270^\circ$ 与顺时针 $90^\circ$ 到达同一点。
取 $-180^\circ<\theta\le180^\circ$ 后，才可用正弦正负区分非平角、非零角的最短转向。
\end{trapbox}
\begin{trapbox}
\textbf{三、反求角度的两个代数检查。}

\textbf{除零检查针对的是 $P=B$，不是 $B=O$。}
若旋转中心恰是原点而 $P\ne B$，反求公式照常可用；若 $P=B$，则 $r^2=0$，
旋转后仍为 $Q=B$，不能除以 $r^2$，也不能由此确定角度。

\textbf{正弦分子顺序不能颠倒。}
正确为 $ut-vs$，错写成 $vs-ut$ 会把角度的正负号翻转。
以 $(u,v)=(1,0)$ 逆时针转 $90^\circ$ 得 $(s,t)=(0,1)$ 作最小自检，
应有 $(ut-vs)/(u^2+v^2)=1$。

\end{trapbox}
